\section{Bounded-degree curves and moving jet kernels}
\label{sec:jet-tools}

The geometric argument will produce subvarieties with bounded degree and
bounded canonical intersection.  We first explain how these bounds
contradict a growing lower bound for the degrees of curves.  We then
record the moving-center results from \cite{LA} that supply the required
subvarieties.  All varieties in this section are over $\C$.

\subsection{The bounded-degree obstruction}

\begin{lemma}[Bounded-degree curves]
\label{lemma:bounded-curves}
For every integer $f\geq1$ there is an integer $m_f>0$ with the
following property.  Let $Y$ be a smooth integral projective variety
of dimension $f$, let $D_0$ be a reduced simple normal crossing divisor
or the zero divisor, and let $L_0$ be a nef integral Cartier divisor
of full nef dimension on $Y$.  Suppose that $P_0$ is a nef rational
Cartier divisor and that $r,C_0,C_1$ are positive real numbers satisfying
\[
 0<P_0^f\leq C_0,\qquad P_0-rL_0\text{ is nef}.
\]
If $A_0=K_Y+D_0+L_0$ is big and
\begin{equation}
\label{j:adjoint-ratio}
 A_0\cdot P_0^{f-1}\leq C_1P_0^f,
\end{equation}
then
\begin{equation}
\label{j:bounded-curve-obstruction}
 r\leq(m_fC_1)^{f-1}C_0.
\end{equation}
The integer $m_f$ depends only on $f$, and not on the variety or
the divisors.  For dimensions $1\leq f\leq n$, one may replace all
$m_f$ by a single integer $m$ depending only on $n$.
\end{lemma}

\begin{proof}
The pair $(Y,D_0)$ is log canonical, its boundary coefficients belong
to $\{0,1\}$, and its nef polarizing divisor $L_0$ has Cartier index
one.  By \cite[Theorem~1.3]{BZ}, an integer $m_f$ depending only on
these data makes $|m_fA_0|$ birational.  Here $A_0$ is integral
Cartier, so no rounding or additional denominator is involved.

Resolve the base ideal by a projective birational morphism
$\mu:\widetilde Y\to Y$, with $\widetilde Y$ smooth, and write
\[
 \mu^*(m_fA_0)=H+F,
\]
where $H$ is the basepoint-free moving part and $F$ is effective.
The morphism defined by $H$ is birational onto its image; in
particular $H$ is nef and big.  In the rest of this proof, $P_0$
and $L_0$ denote their pullbacks to $\widetilde Y$.  Effectivity
of $F$ and nefness of $P_0$ give
\begin{equation}
\label{j:moving-adjoint-degree}
 H\cdot P_0^{f-1}\leq m_fC_1P_0^f.
\end{equation}

For $f\geq2$, put $s_i=P_0^{f-i}\cdot H^i$ for $0\leq i\leq f$.
Both divisors are nef and big, so the $s_i$ are positive.  The
Khovanskii--Teissier inequalities
\cite[Corollary~1.6.3(i) and Example~1.6.4]{LazarsfeldPAGI}
$s_i^2\geq s_{i-1}s_{i+1}$ show
that the ratios $s_i/s_{i-1}$ decrease.  Hence
\begin{equation}
\label{j:mixed-degree}
 P_0\cdot H^{f-1}
 \leq P_0^f\left(\frac{H\cdot P_0^{f-1}}{P_0^f}\right)^{f-1}
 \leq(m_fC_1)^{f-1}C_0.
\end{equation}
Choose a very general point $y\in\widetilde Y$ over the
isomorphism opens of both $\mu$ and the morphism defined by $H$.
We also choose $\mu(y)$ in the very general locus
that tests full nef dimension of $L_0$.  Pull back $f-1$ general
hyperplanes through the image of $y$ under the morphism defined
by $H$.  This subsystem has only the base point $y$, so the
general members meet properly and define an effective curve cycle
of class $H^{f-1}$; the component $C$ through $y$ is a curve
because their local equations are independent at $y$.
Nefness gives
\[
 P_0\cdot C\leq P_0\cdot H^{f-1}.
\]
The curve $C$ is not contracted by $\mu$, and its image passes
through the chosen very general point.  Full nef dimension and
integrality therefore give $L_0\cdot C\geq1$.  Since
$P_0-rL_0$ is nef, we obtain $P_0\cdot C\geq r$, which proves
\eqref{j:bounded-curve-obstruction}.

If $f=1$, take $C=Y$ directly.  Then
$P_0\cdot Y\geq rL_0\cdot Y\geq r$, and the required bound is
$r\leq C_0$.  Finally, \cite[Theorem~1.3]{BZ} applies to every
multiple of its prescribed integer.  A common multiple of the
finitely many integers $m_1,\ldots,m_n$ therefore gives the last
assertion.
\end{proof}

In particular, the hypotheses of Lemma~\ref{lemma:bounded-curves}
cannot hold along a sequence with $r\to\infty$ and with $C_0,C_1$
bounded, even if the varieties and divisors vary.  The following
observation explains how full nef dimension and pseudo-effectivity
of the canonical divisor pass to the moving subvarieties used below.

\begin{lemma}[Very general images]
\label{j:very-general-images}
Let $X$ be a normal integral projective variety that is not uniruled,
and let $L$ be a nef Cartier divisor of full nef dimension on $X$.
There is a complement $X^\circ$ of a countable union of proper
closed subsets of $X$ with the following property.  If $Y$ is a
smooth integral projective variety and $g:Y\to X$ is generically
finite onto a positive-dimensional image meeting $X^\circ$, then
$g^*L$ has full nef dimension, $Y$ is not uniruled, and $K_Y$ is
pseudo-effective.
\end{lemma}

\begin{proof}
Choose $X^\circ$ so that every curve through any of its points has
positive $L$-degree and no rational curve meets $X^\circ$.  The first
condition is the very general characterization of full nef
dimension.  For the second, rational curves are parametrized by
countably many algebraic families; the closure of the evaluation
image of each such family is proper because $X$ is not uniruled.

Since $g(Y)$ meets $X^\circ$, it is contained in none of the closed
sets excluded in its definition.  A very general point of $Y$ thus
maps into $X^\circ$ and lies in the quasi-finite locus of $g$.
Every curve through such a point has a curve as its image, so its
$g^*L$-degree is positive.  This proves full nef dimension.  A
covering family of rational curves on $Y$ would similarly give a
nonconstant rational curve meeting $X^\circ$, a contradiction.
The pseudo-effectivity of $K_Y$ now follows from \cite{BDPP}.
\end{proof}

\subsection{Subadjunction and high-multiplicity components}

We use the following five results from the section
\emph{Tools for moving base components} of \cite{LA}.  Their
hypotheses distinguish generic information along a center from
global information about the ambient variety.  Intersections on
an integral subvariety are computed after pullback to its
normalization and a smooth projective resolution.  A base ideal
is \emph{isolated at the generic point} of a subvariety when its
localization there is primary for the maximal ideal.

\begin{lemma}[Numerical generic subadjunction, \cite{LA}, Lemma~7.1]
\label{j:subadjunction}
Let $Z$ be a projective $\Q$-factorial klt variety, let
$\Theta\geq0$ be a rational divisor, and let $V\subset Z$ be
integral of dimension $f>0$.  Suppose that $(Z,\Theta)$ is log
canonical at the generic point of $V$ and has a divisor of log
discrepancy zero with center $V$.  If $P$ is a nef rational Cartier
divisor on $V$ and $\rho:V^*\to V$ is a smooth projective resolution
factoring through its normalization, then
\[
 K_{V^*}\cdot(\rho^*P)^{f-1}
 \leq (K_Z+\Theta)|_V\cdot P^{f-1}.
\]
\end{lemma}

Only generic log canonicity is required.  The proof in \cite{LA}
uses the dlt modification of \cite[Theorem~2.10]{FH}, whose
truncated boundary leaves an effective surplus supported over
the non-lc locus.  On a minimal stratum over $V$, adjunction
gives a pair that is klt over the generic point of its
connected-fiber base.  This suffices for the klt-trivial
canonical bundle formula and its b-nef moduli divisor
\cite[Definition~3.1 and Theorem~3.5]{FG}.  The effective
discriminant, the moduli divisor, and finite ramification have
nonnegative intersections with the pulled-back nef class.
The comparison uses canonical Weil cycles, so it requires no
$\Q$-Gorenstein assumption on the normalization of $V$.

\begin{lemma}[A base component of high multiplicity, \cite{LA}, Lemma~7.2]
\label{j:base-component}
Let $Z$ be a projective $\Q$-factorial klt variety of dimension
$d$, let $R$ be a nef rational Cartier divisor, and let $k>0$
be an integer such that $kR$ is Cartier.  Let
$0\ne\mathcal H\subset H^0(Z,\OO_Z(kR))$ be a linear subspace
with proper base locus and base ideal $\mathfrak b$.  Suppose that
$V$ is an integral base-locus component of dimension $f<d$,
isolated at its generic point $\eta_V$, and that
$\eta_V\in Z_{\mathrm{sm}}$.  Let
$\mathfrak m$ be the maximal ideal of $\OO_{Z,\eta_V}$.
If $\mathfrak b\OO_{Z,\eta_V}\subset\mathfrak m^h$ for an
integer $h>0$, there are a rational number $0<c\leq(d-f)k/h$ and an
effective rational divisor $\Theta\sim_{\Q}cR$ such that
$(Z,\Theta)$ is log canonical at $\eta_V$ and has a log canonical
place centered on $V$.  The degree and canonical bounds are
\begin{equation}
\label{j:base-bounds}
 \begin{split}
 R^f\cdot V&\leq(k/h)^{d-f}R^d,\\
 K_{V^*}\cdot R^{f-1}
   &\leq(K_Z+cR)|_V\cdot R^{f-1}\qquad(f>0),
 \end{split}
\end{equation}
where $V^*$ is a smooth projective resolution.  The degree bound
also holds for $f=0$.
\end{lemma}

The degree estimate is the local multiplicity bound $h^{d-f}$
combined with successive general cuts by the given linear series.
The boundary is obtained by averaging general members with total
coefficient equal to the generic log canonical threshold of
$\mathfrak b$.  Thus it is an actual divisor to which
Lemma~\ref{j:subadjunction} applies.

\subsection{Full jet kernels and restriction to a fiber}

\begin{lemma}[Moving multiplicity, \cite{LA}, Lemma~7.3]
\label{j:moving-kernel}
Let $Z$ be an integral projective variety of dimension $d$, let
$R$ be a nef, big, semiample rational Cartier divisor, and fix
an integer $k>0$ such that $kR$ is Cartier.  Choose positive real
numbers $\tau_0<\cdots<\tau_d$ with consecutive gap $\delta>0$.
For $x\in Z_{\mathrm{sm}}$, put
\[
 \mathcal H_i(x)=H^0\bigl(Z,\OO_Z(kR)
                  \otimes\mathfrak m_x^{\lceil k\tau_i\rceil}\bigr).
\]
Suppose that for very general $x$ all these spaces are nonzero and
the base locus of $\mathcal H_0(x)$ has a positive-dimensional
component through $x$.  If $k\delta\geq4$, then, after an
algebraic parameter extension and restriction to nonempty opens,
there are an integral parameter variety $T$ with a dominant
marked-point map $x:T\to Z_{\mathrm{sm}}$, a fixed index
$i\in\{0,\ldots,d-1\}$, and a family of integral subvarieties
$V_t$ containing $x(t)$ and common to the base loci of
$\mathcal H_i(x(t))$ and $\mathcal H_{i+1}(x(t))$.  In particular,
the incidence
\[
 \Gamma=\{(t,z):z\in V_t\}\subset T\times Z
\]
dominates $Z$.

For general $t$, the dimension $f$ of $V_t$ lies in $[1,d-1]$,
$V_t$ is an isolated component of the higher base locus at its
generic point, and the entire higher-kernel base ideal is
contained there in the ordinary power $\mathcal I_{V_t}^h$, where
\[
 h=\lceil k\delta/2\rceil,\qquad k/h\leq2/\delta.
\]
The incidence base ideal is likewise contained in
$\mathcal I_\Gamma^h$ on a dense open of $\Gamma$.
These containments are taken in the smooth ambient open.
If $Z$ is also $\Q$-factorial and klt, then
\begin{equation}
\label{j:moving-bounds}
 \begin{split}
 R^f\cdot V_t&\leq(2/\delta)^{d-f}R^d,\\
 K_{V_t^*}\cdot R^{f-1}
   &\leq(K_Z+cR)|_{V_t}\cdot R^{f-1},
       \qquad 0<c\leq2(d-f)/\delta.
 \end{split}
\end{equation}
The top intersection on $V_t$ is positive when $V_t$ meets the
open on which the big semiample contraction is an isomorphism.
\end{lemma}

The use of the full kernels matters: differentiation in the
parameter lowers vanishing at the moving mark and still gives a
global section in the lower kernel.  Incidence dominance makes
these parameter directions span the normal directions.  The
vanishing of all derivatives of order less than $h$ then gives
the ordinary-power containment, which can be restricted to
general fibers as follows.

\begin{lemma}[Restriction of an isolated base ideal, \cite{LA}, Lemma~7.4]
\label{j:fiber-restriction}
In the setting of Lemma~\ref{j:moving-kernel}, let $g:Z\to B$ be
a morphism.  There is a dense open of the incidence with the
following property.  Choose $(t,z)$ in this open, put $b=g(z)$,
and suppose that $V_t$ is smooth over a smooth open of its actual
image at $z$.  Let $F$ be the reduced component through $z$ of
the scheme fiber of $V_t\to g(V_t)$.  Near $z$ in $Z_b$, the
restricted higher-kernel base ideal has support exactly $F$ and
is contained in the ordinary power $\mathcal I_F^h$.

If $F$ is proper in an integral ambient fiber component containing
it, the restricted series is nonzero there and $F$ is an isolated
base component at its generic point.  Applying
Lemma~\ref{j:base-component} on that ambient component additionally
requires its projectivity, $\Q$-factoriality, klt singularities,
and smoothness at the generic point of $F$.
\end{lemma}

Here the fiber is chosen through a general incidence point.  The
statement does not assert the same conclusion for an arbitrary
special fiber.  Smoothness over the actual image identifies the
restricted ideal of $V_t$ with the reduced ideal of $F$ locally;
this is what preserves ordinary powers.

The $\Q$-factoriality needed for the base-component estimates can
also be supplied by a small projective $\Q$-factorialization
$\sigma:Z'\to Z$ of a klt ambient variety.  This passage is made
only for components that meet the isomorphism open of $\sigma$.
Their strict transforms retain the generic base-ideal and
multiplicity conditions for the pulled-back series, and the
testing class $\sigma^*R$ is nef.  Since
$K_{Z'}=\sigma^*K_Z$, projection formula gives the same degree
and canonical-intersection bounds on a common resolution.
The pullback testing class need not be ample; nefness is the
property used in these estimates.  Passing to higher resolutions
does not change the canonical intersections, because their
additional exceptional divisors have zero intersection with
the pulled-back testing class to the required power.

\subsection{Finite covers with a fixed branch complement}

A bound on the degree of a branched cover does not make the
possible covers finite.  The final companion input supplies the
needed fixed branch complement from a family whose branch
divisors do not sweep the base.

\begin{lemma}[Finite covers after excluding sweeping branch divisors,
              \cite{LA}, Lemma~7.5]
\label{j:finite-covers}
Let $Y$ be a normal integral projective variety, let
$\mathcal V\to T$ be a smooth projective family with integral
fibers over an integral parameter variety, and let
$g:\mathcal V\to Y$ be a morphism.  Suppose that the general
fiber maps $g_t:V_t\to Y$ are dominant and generically finite of
degree at most a fixed integer $e$.

After a finite parameter extension and shrinking, spread the
geometric generic components of the zero divisor of the relative
Jacobian over $Y_{\mathrm{sm}}$ to divisors
$\mathcal R_1,\ldots,\mathcal R_N$ in $\mathcal V$.
Retain this notation after the base change.  Let $I$ consist of
the indices $j$ for which
$\overline{g_t((\mathcal R_j)_t)}$ is a divisor in $Y$ for general
$t$.  Assume that
\[
 \overline{g(\mathcal R_j)}\ne Y\qquad(j\in I).
\]
After shrinking again, there is a smooth nonempty open
$Y^\circ\subset Y$ over which every general finite normal Stein
cover of $g_t$ is finite \'etale.  Only finitely many such normal
covers of $Y$ occur, up to isomorphism over $Y$.
The conclusion holds simultaneously for finitely many
projections; the open and the finite lists may depend on all
fixed parameters of the family.
\end{lemma}

To see the role of the hypothesis, remove $Y_{\mathrm{sing}}$
and the finitely many proper image closures
$\overline{g(\mathcal R_j)}$ for $j\in I$.  Every branch prime
of a general Stein cover is accounted for by one of these
relative Jacobian components.  Purity therefore makes the
restricted covers \'etale.  The fundamental group of the
resulting smooth complex variety is finitely generated, so it
has only finitely many covers of degree at most $e$.  Finally,
a finite normal cover of $Y$ is determined by its function
field, hence by its restriction to this fixed open.
